\section{Smooth contact rigidity from the complete boundary law}
\label{sec:v68-smooth-contact}

The boundary law determines functions on a real collar, not merely their
Taylor coefficients.  This distinction removes the analytic hypothesis
from local contact identification.  After the finite jets have been
identified, the actual stationary envelope separates the remaining
functions: its initial evaluation is invertible, whereas every later
evaluation takes place strictly nearer the contact.  A weighted real
norm makes the latter part a contraction.  The argument uses neither
convergence of Taylor series nor a holomorphic extension.

Throughout this section the polygon, the physical phase labels and the
graph sensors are those of Section~\ref{sec:v65-periodic-inverse}.
We write $\mathcal S(F)=(S_b^-)_{b=0}^{r-1}$ for the actual future
actions, with the same anchored gauges.  Contact identification below
means equality on a possibly smaller real neighborhood of each marked
contact.  Continuation to unvisited portions of a smooth obstacle is
not part of that conclusion.

\subsection{A common real collar and successive visits}

\begin{lemma}[Uniform real visits]\label{lem:v68-real-visits}
Let $F^t$, $0\le t\le1$, be a smooth compact family of anchored
$C^\infty$ graph tuples with positive quadratic jets, in the fixed
marked coordinates.  There are $R>0$ and $a<1$ such that the actual
future half-lines are defined for $|u|\le R$ and all $t$, and
\begin{equation}\label{eq:v68-real-contraction}
 x^t_{b,j}(u)=X^t_{b+j-1}\circ\cdots\circ X^t_b(u),\qquad
 |x^t_{b,j}(u)|\le a^j|u|,\qquad j\ge1,
\end{equation}
where $X^t_i(u)=x^t_{i,1}(u)$.  On this collar the real envelope
weights of \eqref{eq:v65-visit-weights} satisfy
\begin{equation}\label{eq:v68-weight-floors}
 \tfrac12\le\alpha_b^t(u)\le\tfrac32,\qquad
 1\le v_{b,j}^t(u)\le3\quad(j\ge1).
\end{equation}
The constants can be chosen uniformly on bounded finite-smoothness
classes with fixed positive curvature bounds and fixed nongrazing
marks.  For each stipulated differentiation order only finitely many
bounds for the graph derivatives are needed.
\end{lemma}
\begin{proof}
The positive Jacobi construction and its weighted implicit equations in
Sections~\ref{sec:v64-periodic-itinerary} and
\ref{sec:v65-periodic-inverse} give the unique decaying stationary
half-line.  Their inverses are uniform on compact positive quadratic
classes.  Applying the construction with the parameter $t$ gives a
common collar and continuous first derivatives of $X_i^t$ there.
By Lemma~\ref{lem:v65-signed-jacobi},
\[
 X_i^t(0)=0,\qquad (X_i^t)'(0)=\sigma_i(t),\qquad
 a_0:=\max_{i,t}|\sigma_i(t)|<1.
\]
Fix $a_0<a<1$ and shrink the collar until
$\sup_{i,t,|u|\le R}|(X_i^t)'(u)|\le a$.  A shifted tail solves
the same decaying stationary problem at the next physical phase.
Uniqueness therefore identifies it with that phase's half-line.
Iteration proves the first identity in
\eqref{eq:v68-real-contraction}; the mean value theorem then proves
the second.  In particular the multiplicative constant is one.
The weaker estimate $C\rho^j|u|$ alone would not suffice for the
large-weight argument below when $C\rho>1$.

Each unit chord is a smooth function of two contact coordinates on a
common neighborhood of the clear reference chord.  There are only
finitely many physical phases.  At the origin the two weights are
$1$ and $2$, respectively.  Uniform continuity, followed by a further
shrink of $R$, gives \eqref{eq:v68-weight-floors} for every incident
pair of chords, hence for every $j$.  Signed tangent choices have not
been changed; these are normal projection weights of actual chords.
For a bounded finite-smoothness class, the same proof uses compactness
of its positive quadratic range and uniform bounds in the weighted
implicit equations.  Bounds through any fixed order follow by
successive differentiation with a strict exponential margin.  No
uniform assertion as curvature or the nongrazing margin vanishes is
used.
\end{proof}

For an integer $m\ge2$ define the weighted real norm
\begin{equation}\label{eq:v68-weighted-norm}
 \|h\|_{m,R}=\max_b\sup_{0<|u|\le R}\frac{|h_b(u)|}{|u|^m}.
\end{equation}
It is finite for a $C^m$ tuple whose derivatives of orders less than
$m$ vanish at zero.  It is a norm on such functions, not a norm of
their formal series.

\begin{proposition}[Separation of the nonlinear remainder]
\label{prop:v68-weighted-separation}
Suppose that two smooth graph tuples $F,G$ have the same contact jets
through order $m-1$.  Apply Lemma~\ref{lem:v68-real-visits} to their
segment $F^t=F+t(G-F)$ on a common collar.  If
\begin{equation}\label{eq:v68-weight-threshold}
 \theta_m:=\frac{6a^m}{1-a^m}<1,
\end{equation}
then
\begin{equation}\label{eq:v68-weighted-inverse}
 \|G-F\|_{m,R}\le
 \frac{2}{1-\theta_m}
       \|\mathcal S(G)-\mathcal S(F)\|_{m,R}.
\end{equation}
The same assertion holds with finite $C^m$ differences whenever the
real half-lines and envelope used here have the indicated regularity.
\end{proposition}
\begin{proof}
Put $h=G-F$ and $g=\mathcal S(G)-\mathcal S(F)$.  Integrate the
actual envelope of Lemma~\ref{lem:v65-envelope} along the segment:
\begin{equation}\label{eq:v68-integrated-envelope}
 g_b(u)=\overline\alpha_b(u)h_b(u)+(Kh)_b(u),\qquad
 \overline\alpha_b(u)=\int_0^1\alpha_b^t(u)\,dt,
\end{equation}
where the linear operator $K$, with this segment fixed, is
\[
 (Kh)_b(u)=\int_0^1\sum_{j\ge1}
       v_{b,j}^t(u)h_{b+j}(x^t_{b,j}(u))\,dt.
\]
The operator is applied to the actual difference, not to a freely
chosen formal variation.  The initial coefficient is independent of
$h_b$ in this displayed identity once the segment is fixed.
For any function in the weighted space,
\[
 |h_{b+j}(x^t_{b,j}(u))|
       \le \|h\|_{m,R}a^{jm}|u|^m.
\]
This summable majorant justifies integration and summation and gives
\[
 \|Kh\|_{m,R}\le\frac{3a^m}{1-a^m}\|h\|_{m,R},\qquad
 \|\overline\alpha^{-1}\|_\infty\le2.
\]
Multiplication by $\overline\alpha^{-1}$ in
\eqref{eq:v68-integrated-envelope} and the triangle inequality prove
\eqref{eq:v68-weighted-inverse}.  Alternatively, the inverse of
$I+\overline\alpha^{-1}K$ is its norm-convergent Neumann series.
The coefficients and composition maps of that operator depend on both
candidates; this proof does not present that pair-dependent inverse as
an algorithm using only the observations.

The finite-jet factorization already proved for actual smooth graphs
shows that $g$ has the same vanishing order as required for this norm.
The terminal contribution in the finite envelope has already been
sent to zero with uniform parameter bounds before the identity is
integrated.  It is not dropped by declaring a formal half-line limit.
\end{proof}

\subsection{Identification of smooth germs}

\begin{theorem}[Smooth marked contact rigidity]
\label{thm:v68-smooth-rigidity}
Fix a clear nongrazing marked polygon with distinct physical contacts.
Two positive-curvature $C^\infty$ contact tuples with the same
phase-resolved two-offset laws \eqref{eq:v65-two-laws} have identical
contact germs in the measured graph coordinates.  No analyticity,
curvature input, flux model or closeness between the candidates is
required.  In fact equality of their actual future action germs at all
physical phases already implies the conclusion.
\end{theorem}
\begin{proof}
The two-offset quotient in Proposition~\ref{prop:v65-two-offset}
recovers the future actions on a common real collar, with the marked
physical momenta restored.  Thus it suffices to treat equality of
$\mathcal S(F)$ and $\mathcal S(G)$ there.
Theorem~\ref{thm:v66-curvature} first identifies the curvatures
without selecting a local inverse branch.  The signed recursion of
Theorem~\ref{thm:v65-cyclic-jets} then identifies every finite contact
jet.  In particular $h=G-F$ is flat at zero at each physical phase.

The segment between these two actual graph representatives has the
same positive quadratic jets.  After shrinking the collar it consists
of admissible local strictly convex arcs with the fixed stationary
polygon.  No interpolation of globally closed obstacles is necessary.
Lemma~\ref{lem:v68-real-visits} supplies one $a<1$ and one common
collar for this segment.  Now choose a finite integer $m\ge3$ so large
that \eqref{eq:v68-weight-threshold} holds.  Flatness and Taylor's
formula make $\|h\|_{m,R}$ finite.  Proposition~\ref{prop:v68-weighted-separation}
with $g=0$ forces this norm to vanish.  Hence $F=G$ on that collar.

Only finitely many jets are used after $a$ has been chosen.  The proof
does not infer equality of smooth functions from their Taylor series:
it uses equality of their full actions once more, in
\eqref{eq:v68-integrated-envelope}.  A flat nonzero contact perturbation
has a nonzero, possibly flat, action response.
\end{proof}

For connected closed analytic obstacle boundaries, fixed placement and
full visitation still give the whole-obstacle conclusion of
Theorem~\ref{thm:v66-global-rigidity}.  For smooth boundaries the
new conclusion concerns the visited contact neighborhoods; arbitrary
unvisited boundary arcs are not determined by local half-line data.
The bounded-holomorphic local inverse of
Theorem~\ref{thm:v65-analytic-inverse} is also retained: equality of
real germs and inversion in a specified analytic Banach norm are
different assertions.

\subsection{Finite-smoothness control from real densities}

The weighted estimate also gives stability of complete real profiles,
rather than only finite jets, under a finite-smoothness prior.  We state
the derivative loss and the range of uniformity explicitly.
Fix a marked polygon, a real interval $I_0=[-R_0,R_0]$, positive
numbers $c_-<c_+$, and a bound $K_0$.  For an integer $m$ let
$\mathcal K_m$ consist of anchored graph tuples satisfying
\begin{equation}\label{eq:v68-smooth-class}
 c_-\le F_i''(0)\le c_+,\qquad
 \max_i\|F_i\|_{C^{m+3}(I_0)}\le K_0.
\end{equation}
We use a sufficiently small common collar of their fixed clear
polygon.  Curvature remains positive there by the $C^3$ bound.
Choose $a<1$ larger than all the one-step linear contraction factors
for the enlarged quadratic box $[c_-/2,2c_+]^r$, and then choose
$m\ge3$ satisfying \eqref{eq:v68-weight-threshold}.  This order is
fixed by the positive marked class; it need not remain bounded in a
weak-hyperbolicity limit.

\begin{lemma}[Finite-order comparison bounds]
\label{lem:v68-finite-comparison}
On a sufficiently small common collar $I=[-R,R]$, the following
bounds hold with constants depending only on the stipulated class.
For $F,G\in\mathcal K_m$, if
$\eta=\|\mathcal S(F)-\mathcal S(G)\|_{C^m(I)}$, then
\begin{equation}\label{eq:v68-jet-alignment-bound}
 \max_{i,\,2\le j<m}|F_i^{(j)}(0)-G_i^{(j)}(0)|\le C\eta.
\end{equation}
If $P$ is an anchored polynomial tuple of degree at most $m-1$,
and the segment $G+tP$ stays in an enlarged bounded positive class,
then
\begin{equation}\label{eq:v68-polynomial-action-bound}
 \|\mathcal S(G+P)-\mathcal S(G)\|_{C^m(I)}
                         \le C\|P\|_{C^{m+3}(I_0)}.
\end{equation}
The enlarged class may be fixed before the small perturbation is chosen.
\end{lemma}
\begin{proof}
For degree two, \eqref{eq:v66-pair-lipschitz} gives a uniform bound
because the positive action Hessians have a common lower bound.
The signed factors depend smoothly on curvature in the enlarged
positive box.  At every subsequent degree subtract the two inversions
\eqref{eq:v65-jet-recursion}.  Their inverse blocks and derivatives
are bounded on that box.  The remainders $\mathcal R_j$ are smooth
functions of the finitely many lower graph jets on bounded boxes with
positive quadratic entries.  Their derivatives there are bounded.
Induction, using the observed action-jet differences, proves
\eqref{eq:v68-jet-alignment-bound}.  This is a two-point comparison;
nonsingularity of an everywhere defined Jacobian is not being used as
a substitute for global injectivity or a global Lipschitz bound.

For \eqref{eq:v68-polynomial-action-bound}, use the weighted
stationary equations along $G+tP$.  On a common collar their
linearized interior operator is a uniformly small perturbation of the
positive Jacobi operator.  It has a uniform inverse on an exponentially
weighted space.  Successively differentiating in the endpoint through
order $m$ puts the same operator on the highest coordinate derivative.
The other terms involve previously bounded derivatives and finitely
many graph derivatives.  A strict margin in the exponential weight
absorbs the fixed polynomial factors in the site index.  With some
$a<\rho<1$, this gives
\[
 |x^t_{b,j}(u)|\le C\rho^j|u|,\qquad
 \sup_{|u|\le R}|\partial_u^k x^t_{b,j}(u)|
                  \le C_k\rho^j\quad(1\le k\le m),
\]
uniformly in $t$.  Chord differentiation likewise bounds
$\|\alpha_b^t\|_{C^m}$ and $\|v_{b,j}^t\|_{C^m}$ uniformly
in $b,j,t$.  The stipulated $C^{m+3}$ graph bound is sufficient for
these finitely differentiated equations and weights.

Because $P(0)=P'(0)=0$, the chain rule gives the useful summable bound
\[
 \|P_{b+j}\circ x^t_{b,j}\|_{C^m([-R,R])}
                  \le C\rho^{2j}\|P\|_{C^m(I_0)}.
\]
For a derivative involving one derivative of $P$, its value is
$O(|x^t_{b,j}|)$ and its coordinate-derivative factor decays as
$\rho^j$.  Terms with at least two derivatives of $P$ have at least
two such coordinate factors.  The undifferentiated term is quadratic
in $x^t_{b,j}$.  These observations prove the bound for every order
at most $m$, including the cases where a polynomial derivative is zero.
Apply this to the actual envelope with variation $P$.  The initial
term has norm at most $C\|P\|_{C^m}$ and the later terms have a
summable $C\rho^{2j}\|P\|_{C^m}$ majorant.  The terminal term
and its fixed derivatives vanish uniformly, as in
Lemma~\ref{lem:v65-envelope}.  Thus the parameter derivative of
the action has $C^m$ norm at most $C\|P\|_{C^{m+3}(I_0)}$.
Integration in $t$ proves the claim.  No optimal derivative count or
same-regularity differentiability theorem for composition operators is
being used.
\end{proof}

\begin{theorem}[Complete real-profile stability]
\label{thm:v68-smooth-stability}
For the class \eqref{eq:v68-smooth-class} and $m$ chosen as above,
there are a common $R>0$, $\eta_0>0$ and $C<\infty$ such that
\begin{equation}\label{eq:v68-real-profile-stability}
 \|F-G\|_{C^0([-R,R])}
       \le C\|\mathcal S(F)-\mathcal S(G)\|_{C^m([-R,R])}
\end{equation}
whenever $F,G\in\mathcal K_m$ and the norm on the right is at most
$\eta_0$.  In particular, on compact positive-density classes with
fixed separated offsets and residual floors,
\begin{equation}\label{eq:v68-law-profile-stability}
 \|F-G\|_{C^0([-R,R])}
 \le C\max_{i,\,d\in\{d_1,d_2\}}
                 \|f^F_{i,d}-f^G_{i,d}\|_{C^m([-R,R]^2)}
\end{equation}
for sufficiently small right side.  These are estimates on realizable
pairs of laws.  They use a finite real-smoothness prior, not an analytic
prior, and do not assert an inverse on arbitrary density tuples.
\end{theorem}
\begin{proof}
Let $\eta$ be the action norm in the statement and set
\[
 P_i(u)=\sum_{j=2}^{m-1}
          \frac{F_i^{(j)}(0)-G_i^{(j)}(0)}{j!}u^j,
 \qquad G^\sharp=G+P.
\]
Lemma~\ref{lem:v68-finite-comparison} gives
$\|P\|_{C^{m+3}(I_0)}\le C\eta$.
For sufficiently small $\eta_0$, both $G+tP$ and the segment
between $F$ and $G^\sharp$ remain in a fixed enlarged positive
class.  Their collars and visit estimates are therefore uniform.
The graph tuples $F,G^\sharp$ have identical jets through order
$m-1$.  By actual finite-jet factorization so do their actions.
Taylor's integral remainder consequently gives
\[
 \|\mathcal S(F)-\mathcal S(G^\sharp)\|_{m,R}
 \le\frac1{m!}
       \|\mathcal S(F)-\mathcal S(G^\sharp)\|_{C^m([-R,R])}
 \le C\eta.
\]
The last inequality uses \eqref{eq:v68-polynomial-action-bound}
and the definition of $\eta$.  Applying
\eqref{eq:v68-weighted-inverse} to this pair, and then removing the
polynomial correction, proves
\[
 \|F-G\|_{C^0}\le R^m\|F-G^\sharp\|_{m,R}+\|P\|_{C^0}
                                                   \le C\eta.
\]
All collars were fixed from the enlarged class before this comparison.
Finally, Proposition~\ref{prop:v65-two-offset} converts the stated
$C^m$ density error to $C^m$ action error on the positive square,
proving \eqref{eq:v68-law-profile-stability}.
\end{proof}

The theorem also proves exact contact identification on each of these
finite-smoothness classes.  Its quantitative conclusion is a $C^0$
profile estimate from a $C^m$ law error with a $C^{m+3}$ graph bound.
It is not an analytic-norm estimate, a derivative-free statistical
rate, or a theorem uniform in the period.  If two realizable candidates
fit common finite-flight density records to accuracy $\varepsilon$
in this same $C^m$ norm, the relative theorem at order $m$ gives
\[
 \|F-G\|_{C^0([-R,R])}\le C(\tau^N+\varepsilon)
\]
whenever the combined error is small and the forward constants are
uniform on the class.  This uses the triangle inequality through the
common records.  It does not manufacture density derivative bounds
from total variation or assign a preparation budget.  The curvature
stopping error in Proposition~\ref{prop:v66-finite-flight} remains a
separate, explicitly propagated computational error.

\subsection{A smooth physical distinction beyond all contact jets}

\begin{proposition}[Flat contacts distinguished by boundary laws]
\label{prop:v68-flat-family}
Let a smooth strictly dispersing table have a clear nongrazing marked
periodic polygon with distinct contact points.  There is a small smooth
family of actual tables, with the same polygon and free cell area,
whose complete contact Taylor jets at every marked contact agree, but
whose two-offset boundary laws do not agree on any sufficiently small
common contact square.  Every future action also has the same Taylor jets throughout the
family.  Thus a formal all-order contact record and the complete
boundary-law record have different geometric information even in the
fixed marked physical class.
\end{proposition}
\begin{proof}
Choose one contact, called $b$, and a smooth nonnegative cutoff $\chi$
supported in its graph chart and equal to one near zero.  Replace its
graph by
\[
 F_b^\varepsilon(u)=F_b^0(u)+\varepsilon\psi(u),\qquad
 \psi(u)=\begin{cases}\chi(u)e^{-1/u^2},&u\ne0,\\0,&u=0.
             \end{cases}
\]
All other contact graphs are unchanged.  The difference is flat but
nonzero on every punctured neighborhood.  For small $\varepsilon$,
strict convexity, embeddedness, disjointness and the clearance margins
persist.  The fixed contact points, tangent frames and curvatures
preserve the original periodic polygon exactly.

To keep the free area fixed, choose a smooth normal bump on a boundary
arc disjoint from every contact collar, with nonzero first area
variation.  Such an arc exists after the finitely many collars are
shrunk.  For a positively oriented obstacle boundary $R(s)$, its area
is $\frac12\int R\mathbin\times R'\,ds$; differentiation and
integration by parts give first variation
$\int\delta R\mathbin\times R'\,ds$.  A nonzero nonnegative
outward normal bump therefore has nonzero area derivative.  The ordinary
implicit function theorem chooses its coefficient smoothly in
$\varepsilon$ so that the obstacle area, hence the free cell area, is
unchanged.  The correction is outside all the relevant charts and
leaves their stationary half-lines unchanged.  This constructs actual
closed obstacles, rather than only formal local coefficients.

All graph jets remain fixed.  The actual finite-jet factorization
therefore fixes every action jet as well.  The exact linear reference
twists, flight lengths, incidence angles and multipliers are also
unchanged.  Nevertheless the envelope on the perturbed contact detects
the difference.  On a smaller collar, \eqref{eq:v68-weight-floors}
and positivity of $\psi$ give, for $\varepsilon>0$ and $u\ne0$,
\[
 S_b^{-,\varepsilon}(u)-S_b^{-,0}(u)
       \ge\frac{\varepsilon}{2}e^{-1/u^2}>0.
\]
The area-correcting bump contributes nothing to this local action.
More precisely, uniformly for a sufficiently small parameter range,
\begin{equation}\label{eq:v68-flat-signal}
 S_b^{-,\varepsilon}(u)-S_b^{-,0}(u)
   =\varepsilon e^{-1/u^2}
       \left\{1+O(|u|)
           +O\!\left(e^{-(a^{-2}-1)/u^2}\right)\right\}.
\end{equation}
Indeed $\alpha_b^t(u)=1+O(|u|)$, while each later perturbed
visit has $|x^t_{b,j}(u)|\le a^j|u|$.  Dividing its flat value
by $e^{-1/u^2}$ and using
$a^{-2j}-1\ge j(a^{-2}-1)$ bounds the entire tail by the convergent
geometric series with ratio $e^{-(a^{-2}-1)/u^2}$.
Visits to unperturbed phases contribute zero.

Equality of the two-offset laws would imply equality of these actions
by the anchored quotient \eqref{eq:v65-offset-quotient}, a
contradiction.  The claim is not equality of the complete fixed-window
count or of the normalized density Taylor coefficients: their
normalizers can change.  It is equality of all formal contact and
action data, with the reference geometry and area fixed, in the presence
of a nonzero physical-law signal.  Nor is it a comparison of full marked
length spectra, which are not the observations used here.
\end{proof}

The flat family explains why the complete-law hypothesis in
Theorem~\ref{thm:v68-smooth-rigidity} is substantive.  It is not a
formal consequence of all-order contact separation.  Exact relative
normalization first retains the real action functions in the physical
law; the contracted visits then identify those functions with actual
boundary profiles, including their flat parts.  The finite-smoothness
estimate records the derivative information needed for stable use of
that same mechanism.
